\subsection{Worked solutions on the ten soils (second scenario)}
The same questions are repeated for different working conditions, to show how the answer moves with them.

\exbox{Ten-soil example 2b: Coarse loam (S01)}{Coarse loam: $K_s=31.2$ cm\,d$^{-1}$ (typical for the class), $\alpha=0.0961$, $n=1.353$. Find $K(h)$ at 100, 300 and 1000 cm suction and the steady upward flux between $h=300$ cm at 0.4 m and $h=100$ cm at 0 m, using Darcy's law $q=K(\Delta h/\Delta z-1)$ with the geometric mean of $K$.}\sol \givenline{$K_s=31.2$ cm\,d$^{-1}$ ; $\alpha=0.0961$ ; $n=1.353$ ; $h=300$ cm ; $h=100$ cm ; $q=K(\Delta h/\Delta z-1)$}\ans $K_r(100)=9.39\times10^{-5}$, $K_r(300)=4.16\times10^{-6}$, $K_r(1000)=1.31\times10^{-7}$ \ans $K=0.00293$, $0.00013$, $4.08\times10^{-6}$ cm\,d$^{-1}$ \ans Geometric mean of the two ends: \[\bar K=\sqrt{0.00293\times0.00013}=0.000616\ \text{cm\,d$^{-1}$}\] \ans Gradient: \[\Delta h/\Delta z=(300-100)/40=5.0\]
so \[q=0.000616\times(5.0-1)=0.00246\ \text{cm\,d$^{-1}$}\] \ans The flux is negligible against the demand of a crop (0.5 cm per day). \finans

\exbox{Ten-soil example 2b: Coarse sand (S02)}{Coarse sand: $K_s=288.0$ cm\,d$^{-1}$ (typical for the class), $\alpha=0.0536$, $n=1.600$. Find $K(h)$ at 100, 300 and 1000 cm suction and the steady upward flux between $h=300$ cm at 0.4 m and $h=100$ cm at 0 m, using Darcy's law $q=K(\Delta h/\Delta z-1)$ with the geometric mean of $K$.}\sol \givenline{$K_s=288.0$ cm\,d$^{-1}$ ; $\alpha=0.0536$ ; $n=1.600$ ; $h=300$ cm ; $h=100$ cm ; $q=K(\Delta h/\Delta z-1)$}\ans $K_r(100)=0.000356$, $K_r(300)=8.28\times10^{-6}$, $K_r(1000)=1.24\times10^{-7}$ \ans $K=0.102$, $0.00239$, $3.58\times10^{-5}$ cm\,d$^{-1}$ \ans Geometric mean of the two ends: \[\bar K=\sqrt{0.102\times0.00239}=0.0156\ \text{cm\,d$^{-1}$}\] \ans Gradient: \[\Delta h/\Delta z=(300-100)/40=5.0\]
so \[q=0.0156\times(5.0-1)=0.0625\ \text{cm\,d$^{-1}$}\] \ans The flux is negligible against the demand of a crop (0.5 cm per day). \finans

\exbox{Ten-soil example 2b: Sand (S03)}{Sand: $K_s=504.0$ cm\,d$^{-1}$ (typical for the class), $\alpha=0.1450$, $n=2.680$. Find $K(h)$ at 100, 300 and 1000 cm suction and the steady upward flux between $h=300$ cm at 0.4 m and $h=100$ cm at 0 m, using Darcy's law $q=K(\Delta h/\Delta z-1)$ with the geometric mean of $K$.}\sol \givenline{$K_s=504.0$ cm\,d$^{-1}$ ; $\alpha=0.1450$ ; $n=2.680$ ; $h=300$ cm ; $h=100$ cm ; $q=K(\Delta h/\Delta z-1)$}\ans $K_r(100)=2.47\times10^{-8}$, $K_r(300)=2.73\times10^{-11}$, $K_r(1000)=1.56\times10^{-14}$ \ans $K=1.25\times10^{-5}$, $1.37\times10^{-8}$, $7.88\times10^{-12}$ cm\,d$^{-1}$ \ans Geometric mean of the two ends: \[\bar K=\sqrt{1.25\times10^{-5}\times1.37\times10^{-8}}=4.14\times10^{-7}\ \text{cm\,d$^{-1}$}\] \ans Gradient: \[\Delta h/\Delta z=(300-100)/40=5.0\]
so \[q=4.14\times10^{-7}\times(5.0-1)=1.66\times10^{-6}\ \text{cm\,d$^{-1}$}\] \ans The flux is negligible against the demand of a crop (0.5 cm per day). \finans

\exbox{Ten-soil example 2b: Loamy sand (S04)}{Loamy sand: $K_s=146.4$ cm\,d$^{-1}$ (typical for the class), $\alpha=0.1240$, $n=2.280$. Find $K(h)$ at 100, 300 and 1000 cm suction and the steady upward flux between $h=300$ cm at 0.4 m and $h=100$ cm at 0 m, using Darcy's law $q=K(\Delta h/\Delta z-1)$ with the geometric mean of $K$.}\sol \givenline{$K_s=146.4$ cm\,d$^{-1}$ ; $\alpha=0.1240$ ; $n=2.280$ ; $h=300$ cm ; $h=100$ cm ; $q=K(\Delta h/\Delta z-1)$}\ans $K_r(100)=6.46\times10^{-7}$, $K_r(300)=2.15\times10^{-9}$, $K_r(1000)=4.1\times10^{-12}$ \ans $K=9.46\times10^{-5}$, $3.14\times10^{-7}$, $6\times10^{-10}$ cm\,d$^{-1}$ \ans Geometric mean of the two ends: \[\bar K=\sqrt{9.46\times10^{-5}\times3.14\times10^{-7}}=5.45\times10^{-6}\ \text{cm\,d$^{-1}$}\] \ans Gradient: \[\Delta h/\Delta z=(300-100)/40=5.0\]
so \[q=5.45\times10^{-6}\times(5.0-1)=2.18\times10^{-5}\ \text{cm\,d$^{-1}$}\] \ans The flux is negligible against the demand of a crop (0.5 cm per day). \finans

\exbox{Ten-soil example 2b: Sandy loam (S05)}{Sandy loam: $K_s=62.4$ cm\,d$^{-1}$ (typical for the class), $\alpha=0.0750$, $n=1.890$. Find $K(h)$ at 100, 300 and 1000 cm suction and the steady upward flux between $h=300$ cm at 0.4 m and $h=100$ cm at 0 m, using Darcy's law $q=K(\Delta h/\Delta z-1)$ with the geometric mean of $K$.}\sol \givenline{$K_s=62.4$ cm\,d$^{-1}$ ; $\alpha=0.0750$ ; $n=1.890$ ; $h=300$ cm ; $h=100$ cm ; $q=K(\Delta h/\Delta z-1)$}\ans $K_r(100)=4.29\times10^{-5}$, $K_r(300)=4.27\times10^{-7}$, $K_r(1000)=2.65\times10^{-9}$ \ans $K=0.00268$, $2.67\times10^{-5}$, $1.65\times10^{-7}$ cm\,d$^{-1}$ \ans Geometric mean of the two ends: \[\bar K=\sqrt{0.00268\times2.67\times10^{-5}}=0.000267\ \text{cm\,d$^{-1}$}\] \ans Gradient: \[\Delta h/\Delta z=(300-100)/40=5.0\]
so \[q=0.000267\times(5.0-1)=0.00107\ \text{cm\,d$^{-1}$}\] \ans The flux is negligible against the demand of a crop (0.5 cm per day). \finans

\exbox{Ten-soil example 2b: Loam (S06)}{Loam: $K_s=31.2$ cm\,d$^{-1}$ (typical for the class), $\alpha=0.0360$, $n=1.560$. Find $K(h)$ at 100, 300 and 1000 cm suction and the steady upward flux between $h=300$ cm at 0.4 m and $h=100$ cm at 0 m, using Darcy's law $q=K(\Delta h/\Delta z-1)$ with the geometric mean of $K$.}\sol \givenline{$K_s=31.2$ cm\,d$^{-1}$ ; $\alpha=0.0360$ ; $n=1.560$ ; $h=300$ cm ; $h=100$ cm ; $q=K(\Delta h/\Delta z-1)$}\ans $K_r(100)=0.00136$, $K_r(300)=3.8\times10^{-5}$, $K_r(1000)=6.55\times10^{-7}$ \ans $K=0.0424$, $0.00119$, $2.04\times10^{-5}$ cm\,d$^{-1}$ \ans Geometric mean of the two ends: \[\bar K=\sqrt{0.0424\times0.00119}=0.00709\ \text{cm\,d$^{-1}$}\] \ans Gradient: \[\Delta h/\Delta z=(300-100)/40=5.0\]
so \[q=0.00709\times(5.0-1)=0.0284\ \text{cm\,d$^{-1}$}\] \ans The flux is negligible against the demand of a crop (0.5 cm per day). \finans

\exbox{Ten-soil example 2b: Silt loam (S07)}{Silt loam: $K_s=16.3$ cm\,d$^{-1}$ (typical for the class), $\alpha=0.0200$, $n=1.410$. Find $K(h)$ at 100, 300 and 1000 cm suction and the steady upward flux between $h=300$ cm at 0.4 m and $h=100$ cm at 0 m, using Darcy's law $q=K(\Delta h/\Delta z-1)$ with the geometric mean of $K$.}\sol \givenline{$K_s=16.3$ cm\,d$^{-1}$ ; $\alpha=0.0200$ ; $n=1.410$ ; $h=300$ cm ; $h=100$ cm ; $q=K(\Delta h/\Delta z-1)$}\ans $K_r(100)=0.00652$, $K_r(300)=0.000335$, $K_r(1000)=9.6\times10^{-6}$ \ans $K=0.106$, $0.00547$, $0.000157$ cm\,d$^{-1}$ \ans Geometric mean of the two ends: \[\bar K=\sqrt{0.106\times0.00547}=0.0241\ \text{cm\,d$^{-1}$}\] \ans Gradient: \[\Delta h/\Delta z=(300-100)/40=5.0\]
so \[q=0.0241\times(5.0-1)=0.0964\ \text{cm\,d$^{-1}$}\] \ans The flux is negligible against the demand of a crop (0.5 cm per day). \finans

\exbox{Ten-soil example 2b: Sandy clay loam (S08)}{Sandy clay loam: $K_s=10.3$ cm\,d$^{-1}$ (typical for the class), $\alpha=0.0590$, $n=1.480$. Find $K(h)$ at 100, 300 and 1000 cm suction and the steady upward flux between $h=300$ cm at 0.4 m and $h=100$ cm at 0 m, using Darcy's law $q=K(\Delta h/\Delta z-1)$ with the geometric mean of $K$.}\sol \givenline{$K_s=10.3$ cm\,d$^{-1}$ ; $\alpha=0.0590$ ; $n=1.480$ ; $h=300$ cm ; $h=100$ cm ; $q=K(\Delta h/\Delta z-1)$}\ans $K_r(100)=0.000324$, $K_r(300)=1.05\times10^{-5}$, $K_r(1000)=2.26\times10^{-7}$ \ans $K=0.00334$, $0.000108$, $2.33\times10^{-6}$ cm\,d$^{-1}$ \ans Geometric mean of the two ends: \[\bar K=\sqrt{0.00334\times0.000108}=0.0006\ \text{cm\,d$^{-1}$}\] \ans Gradient: \[\Delta h/\Delta z=(300-100)/40=5.0\]
so \[q=0.0006\times(5.0-1)=0.0024\ \text{cm\,d$^{-1}$}\] \ans The flux is negligible against the demand of a crop (0.5 cm per day). \finans

\exbox{Ten-soil example 2b: Clay loam (S09)}{Clay loam: $K_s=5.5$ cm\,d$^{-1}$ (typical for the class), $\alpha=0.0190$, $n=1.310$. Find $K(h)$ at 100, 300 and 1000 cm suction and the steady upward flux between $h=300$ cm at 0.4 m and $h=100$ cm at 0 m, using Darcy's law $q=K(\Delta h/\Delta z-1)$ with the geometric mean of $K$.}\sol \givenline{$K_s=5.5$ cm\,d$^{-1}$ ; $\alpha=0.0190$ ; $n=1.310$ ; $h=300$ cm ; $h=100$ cm ; $q=K(\Delta h/\Delta z-1)$}\ans $K_r(100)=0.00574$, $K_r(300)=0.000392$, $K_r(1000)=1.54\times10^{-5}$ \ans $K=0.0317$, $0.00216$, $8.5\times10^{-5}$ cm\,d$^{-1}$ \ans Geometric mean of the two ends: \[\bar K=\sqrt{0.0317\times0.00216}=0.00828\ \text{cm\,d$^{-1}$}\] \ans Gradient: \[\Delta h/\Delta z=(300-100)/40=5.0\]
so \[q=0.00828\times(5.0-1)=0.0331\ \text{cm\,d$^{-1}$}\] \ans The flux is negligible against the demand of a crop (0.5 cm per day). \finans

\exbox{Ten-soil example 2b: Silty clay loam (S10)}{Silty clay loam: $K_s=3.6$ cm\,d$^{-1}$ (typical for the class), $\alpha=0.0100$, $n=1.230$. Find $K(h)$ at 100, 300 and 1000 cm suction and the steady upward flux between $h=300$ cm at 0.4 m and $h=100$ cm at 0 m, using Darcy's law $q=K(\Delta h/\Delta z-1)$ with the geometric mean of $K$.}\sol \givenline{$K_s=3.6$ cm\,d$^{-1}$ ; $\alpha=0.0100$ ; $n=1.230$ ; $h=300$ cm ; $h=100$ cm ; $q=K(\Delta h/\Delta z-1)$}\ans $K_r(100)=0.0139$, $K_r(300)=0.00153$, $K_r(1000)=8.64\times10^{-5}$ \ans $K=0.0499$, $0.00551$, $0.000311$ cm\,d$^{-1}$ \ans Geometric mean of the two ends: \[\bar K=\sqrt{0.0499\times0.00551}=0.0166\ \text{cm\,d$^{-1}$}\] \ans Gradient: \[\Delta h/\Delta z=(300-100)/40=5.0\]
so \[q=0.0166\times(5.0-1)=0.0663\ \text{cm\,d$^{-1}$}\] \ans The flux is negligible against the demand of a crop (0.5 cm per day). \finans

